% Complete reviewed Unit 2; editorial master: chapters/ch02/student.md
% Markdown SHA-256: 7ac884d55c143f1e5ce6d5eead3b09ec11ea9121cb5f23f4c96e0a9addb0dea2
% One-use initial assembly followed by complete content and page review.
\newcommand{\runningtitle}{Unit 2. Pronouns and weak verbs}
% Layout inherited from the reviewed previous course; set a running title per unit.
% This is a small hand-maintained TeX layout, not a conversion pipeline.
\RequirePackage{fix-cm}
\documentclass[11pt,twoside]{article}
\usepackage{fontspec}
\setmainfont{Cambria}[
  Renderer=HarfBuzz,
  BoldFont=cambriab.ttf,
  ItalicFont=cambriai.ttf,
  BoldItalicFont=cambriaz.ttf
]
\setmonofont{lmmono10-regular.otf}[BoldFont=lmmonolt10-bold.otf]
\usepackage[paperwidth=165mm,paperheight=241mm,inner=21mm,outer=19mm,
  top=22mm,bottom=20mm,headheight=14pt,headsep=7mm,footskip=10mm]{geometry}
\usepackage{xcolor,array,booktabs,longtable,tabularx}
\usepackage{enumitem,fancyhdr}
\usepackage[unicode,hidelinks]{hyperref}
\definecolor{ink}{HTML}{000000}
\definecolor{muted}{HTML}{4C5558}
\setlength{\parindent}{0pt}
\setlength{\parskip}{3.5pt}
\AtBeginDocument{\fontsize{11}{14.3}\selectfont}
\setlength{\tabcolsep}{4pt}
\renewcommand{\arraystretch}{1.10}
\setlist[itemize]{leftmargin=1.3em,label=\textbullet,itemsep=2pt,topsep=3pt,parsep=0pt}
\setlist[enumerate]{leftmargin=1.8em,itemsep=3pt,topsep=4pt,parsep=0pt}
\widowpenalty=10000
\clubpenalty=10000
\emergencystretch=1em
\raggedbottom
\pagestyle{fancy}
\fancyhf{}
\providecommand{\runningtitle}{Old English}
\fancyhead[LE,RO]{\small\itshape \runningtitle}
\fancyhead[LO,RE]{\small Old English}
\fancyfoot[LE,RO]{\small\thepage}
\renewcommand{\headrulewidth}{0.3pt}
\newcommand{\booktitle}[2]{
  {\fontsize{9}{11}\selectfont\color{muted}#1\par}
  \vspace{3pt}
  {\fontsize{20}{24}\selectfont\color{ink}#2\par}
  \vspace{6pt}
}
\newcommand{\source}[1]{\par{\footnotesize\color{muted}#1\par}}
\newcommand{\readingtext}[1]{
  \par\addvspace{5pt}
  \begingroup\leftskip=1.5em
    \fontsize{12}{16}\selectfont #1\par
  \endgroup\addvspace{5pt}
}
\makeatletter
\renewcommand{\section}{\@startsection{section}{1}{\z@}
  {-2.4ex \@plus -0.6ex \@minus -0.2ex}{0.8ex \@plus 0.2ex}
  {\normalfont\fontsize{12}{15}\selectfont\bfseries\color{ink}}}
\renewcommand{\subsection}{\@startsection{subsection}{2}{\z@}
  {-1.8ex \@plus -0.5ex \@minus -0.2ex}{0.5ex \@plus 0.2ex}
  {\normalfont\fontsize{11}{14}\selectfont\bfseries\color{ink}}}
\makeatother

\hyphenpenalty=10000
\exhyphenpenalty=10000
\emergencystretch=3em
\newcommand{\itempara}[2]{\par\addvspace{3pt}\noindent\hangindent=1.8em\hangafter=1\makebox[1.8em][l]{#1.}#2\par}
\newcommand{\reserve}[1]{\par\begingroup\dimen0=\pagegoal\advance\dimen0 by-\pagetotal\ifdim\dimen0<#1\newpage\fi\endgroup}
\begin{document}
\raggedright
\booktitle{UNIT 2}{Personal pronouns and weak indicative verbs}

% mdparagraph:P001:begin
This unit introduces personal pronouns and the present and preterite indicative of weak verbs. Learn the finite forms in the tables, together with their stem changes. The six first- and second-person genitives and the negative verb \textbf{nabban} are for recognition. Infinitive and participle citations are available for lookup; their uses in sentences come later.
% mdparagraph:P001:end

\reserve{7\baselineskip}
% mdparagraph:P002:begin
\section*{1. Personal pronouns}
% mdparagraph:P002:end

% mdparagraph:P003:begin
The first person identifies the speaker; the second identifies the addressee; the third refers to someone or something else. Besides singular and plural, the first and second persons have a \textbf{dual}, explicitly referring to two. A dual subject takes the corresponding \textbf{plural verb form}; there is no separate dual verb ending. Use the dual when an exercise requests “we two” or “you two”. Ordinary plural forms are not restricted by a rule that they must always exclude two. [1]
% mdparagraph:P003:end

% teaching: NEW-U02-PAR-ic NEW-U02-PAR-thu NEW-U02-PAR-he NEW-U02-VOC-ic NEW-U02-VOC-thu NEW-U02-VOC-he NEW-U02-GRAM-personal-pronouns NEW-U02-GRAM-personal-genitive-citation NEW-U02-GRAM-dual-agreement
% mdtable:T01:begin
\par\begingroup
\fontsize{10.5}{13.2}\selectfont
\setlength{\LTleft}{0pt}\setlength{\LTright}{\fill}
\noindent
\begin{tabular}{@{}>{\raggedright\arraybackslash}p{28mm}>{\raggedright\arraybackslash}p{22mm}>{\raggedright\arraybackslash}p{20.5mm}>{\raggedright\arraybackslash}p{21mm}>{\raggedright\arraybackslash}p{20.5mm}@{}}
\toprule
\textbf{Person and number} & \textbf{Nominative} & \textbf{Genitive} & \textbf{Dative} & \textbf{Accusative} \\
\midrule
1 singular: I & ic & mīn & mē & mē \\
1 dual: we two & wit & uncer & unc & unc \\
1 plural: we & wē & ūre & ūs & ūs \\
2 singular: you & ðū & ðīn & ðē & ðē \\
2 dual: you two & git & incer & inc & inc \\
2 plural: you & gē & ēower & ēow & ēow \\
3 masculine singular & hē & his & him & hine \\
3 neuter singular & hit & his & him & hit \\
3 feminine singular & hēo & hiere & hiere & hīe \\
3 plural: all genders & hīe & hiera & him & hīe \\
\bottomrule
\end{tabular}
\par\endgroup

% mdtable:T01:end
% mdparagraph:P004:begin
Recognize \textbf{mīn, uncer, ūre, ðīn, incer, ēower} as genitives of the persons and numbers shown. Their corresponding possessive adjectives have endings not yet taught. Do not use this table alone to form attributive “my”, “your” or “our” phrases. Learn the remaining cells for use in the constructions below.
% mdparagraph:P004:end

% teaching: NEW-U02-GRAM-pronoun-reference
% mdparagraph:P005:begin
The nominative is the subject form. The accusative can be a direct object; the dative can name a recipient. Their spelling often coincides: \textbf{mē} and \textbf{ðē} each serve two cases. \textbf{Him} can be singular or plural; \textbf{hīe} can be feminine singular accusative or plural nominative/\allowbreak{}accusative. The sentence and its context settle the analysis.
% mdparagraph:P005:end

% mdparagraph:P006:begin
Third-person pronouns follow the grammatical gender and number of their referents. English “it” therefore need not correspond to \textbf{hit}: a feminine noun can be resumed by \textbf{hēo}, a masculine one by \textbf{hē}. Case still follows the pronoun’s role in its own clause.
% mdparagraph:P006:end

% mdparagraph:P007:begin
The table selects one spelling for each cell. Cook also has \textbf{mec, ðec, ūsic, ēowic} as accusatives, \textbf{ūser} as first-person plural genitive, \textbf{hiene} beside \textbf{hine}, \textbf{hire} beside \textbf{hiere}, and \textbf{hira/\allowbreak{}heora} beside \textbf{hiera}. Other third-person spellings include \textbf{hīo/\allowbreak{}hī} for feminine nominative, \textbf{hī/\allowbreak{}hēo} for feminine accusative and plural nominative/\allowbreak{}accusative. These are reference alternatives, not extra recall requirements. Preserve the quantities in the selected forms: \textbf{hīe}, but \textbf{hine, hiere, hiera}.
% mdparagraph:P007:end

% teaching: NEW-U02-GRAM-third-person-possessors
% mdparagraph:P008:begin
Third-person genitives can be used as invariant possessors: \textbf{his}, “his/\allowbreak{}its”; \textbf{hiere}, “her”; \textbf{hiera}, “their”. Choose the form by the possessor. It does not change to agree with the possessed noun, which retains its own case and number. Thus the singular or plural of a possessed noun does not change \textbf{his}; a dative possessed noun does not turn \textbf{hiera} into \textbf{him}. These genitives precede the noun without an added article in our model phrases. [2]
% mdparagraph:P008:end

\reserve{7\baselineskip}
% mdparagraph:P009:begin
\section*{2. Reading the verb tables}
% mdparagraph:P009:end

% teaching: NEW-U02-GRAM-weak-indicative-scope
% mdparagraph:P010:begin
Weak verbs form their preterite with a dental element, \textbf{d} or \textbf{t}, sometimes accompanied by a stem change. The \textbf{indicative} states an action or state; the \textbf{preterite} is the past tense. The three classes below must be learned with the verb: the infinitive ending alone does not identify the class.
% mdparagraph:P010:end

% mdparagraph:P011:begin
Every table gives the complete singular and plural present and preterite indicative. \textbf{1/\allowbreak{}3 sg.} means both first and third singular; \textbf{all pl.} means first, second and third plural. The plural column also supplies the form for dual subjects. Thus seven displayed forms account for all twelve person–number–tense cells.
% mdparagraph:P011:end

% mdparagraph:P012:begin
In the selected preterite, the first and third singular coincide. Add \textbf{-st} to that form for the second singular; replace its final \textbf{-e} with \textbf{-on} for all three plural persons. The present singular distinguishes three persons, while all three plural persons coincide. Observe the actual stems and contractions; do not simply attach endings to an infinitive. [3]
% mdparagraph:P012:end

\reserve{7\baselineskip}
% mdparagraph:P013:begin
\section*{3. Class I: long stems and dental changes}
% mdparagraph:P013:end

% teaching: NEW-U02-GRAM-weak-I-long
% mdparagraph:P014:begin
The basic present endings are \textbf{-e, -(e)st, -(e)ð, -að}. In the longer stems, loss of the unstressed \textbf{e} in the second and third singular brings consonants together. The resulting forms must be learned with the preterite. Compare all ten patterns in the following tables. [4]
% mdparagraph:P014:end

% teaching: NEW-U02-PAR-hieran NEW-U02-PAR-gefyllan NEW-U02-PAR-cyssan NEW-U02-PAR-settan NEW-U02-PAR-sendan
% mdtable:T02:begin
\par\begin{minipage}{\linewidth}\raggedright
\fontsize{10.5}{13.2}\selectfont
% panel:Present indicative
\textbf{Present indicative}\par\smallskip
\setlength{\LTleft}{0pt}\setlength{\LTright}{\fill}
\noindent
\begin{tabular}{@{}>{\raggedright\arraybackslash}p{25.5mm}>{\raggedright\arraybackslash}p{21mm}>{\raggedright\arraybackslash}p{22mm}>{\raggedright\arraybackslash}p{22mm}>{\raggedright\arraybackslash}p{22mm}@{}}
\toprule
\textbf{Verb} & \textbf{1 sg.} & \textbf{2 sg.} & \textbf{3 sg.} & \textbf{All pl.} \\
\midrule
hīeran & hīere & hīerst & hīerð & hīerað \\
gefyllan & gefylle & gefylst & gefylð & gefyllað \\
cyssan & cysse & cyst & cyst & cyssað \\
sęttan & sętte & sętst & sętt & sęttað \\
sęndan & sęnde & sęntst & sęnt & sęndað \\
\bottomrule
\end{tabular}
\par\medskip
% panel:Preterite indicative
\textbf{Preterite indicative}\par\smallskip
\setlength{\LTleft}{0pt}\setlength{\LTright}{\fill}
\noindent
\begin{tabular}{@{}>{\raggedright\arraybackslash}p{28mm}>{\raggedright\arraybackslash}p{28mm}>{\raggedright\arraybackslash}p{33mm}>{\raggedright\arraybackslash}p{26mm}@{}}
\toprule
\textbf{Verb} & \textbf{1/\allowbreak{}3 sg.} & \textbf{2 sg.} & \textbf{All pl.} \\
\midrule
hīeran & hīerde & hīerdest & hīerdon \\
gefyllan & gefylde & gefyldest & gefyldon \\
cyssan & cyste & cystest & cyston \\
sęttan & sętte & sęttest & sętton \\
sęndan & sęnde & sęndest & sęndon \\
\bottomrule
\end{tabular}
\par\medskip
\end{minipage}\par

% mdtable:T02:end
% mdparagraph:P015:begin
\textbf{Hīeran} has ordinary \textbf{-de}. In \textbf{gefyllan}, doubled \textbf{ll} becomes single before another consonant: \textbf{gefylst, gefylð, gefylde}. The \textbf{ss/\allowbreak{}ff/\allowbreak{}pp} family has \textbf{-te} with simplification, here \textbf{cyssan → cyste}. Its contracted present second and third singular are both \textbf{cyst}. \textbf{Sęttan} retains \textbf{tt} in \textbf{sętte}; \textbf{sęndan} has \textbf{sęnde}, without an added written dental. First singular present and first/\allowbreak{}third singular preterite can consequently look alike.
% mdparagraph:P015:end

% teaching: NEW-U02-PAR-laedan NEW-U02-PAR-iecan NEW-U02-PAR-gerestan NEW-U02-PAR-metan NEW-U02-PAR-gierwan
% mdtable:T03:begin
\par\begin{minipage}{\linewidth}\raggedright
\fontsize{10.5}{13.2}\selectfont
% panel:Present indicative
\textbf{Present indicative}\par\smallskip
\setlength{\LTleft}{0pt}\setlength{\LTright}{\fill}
\noindent
\begin{tabular}{@{}>{\raggedright\arraybackslash}p{25.5mm}>{\raggedright\arraybackslash}p{21mm}>{\raggedright\arraybackslash}p{22mm}>{\raggedright\arraybackslash}p{22mm}>{\raggedright\arraybackslash}p{22mm}@{}}
\toprule
\textbf{Verb} & \textbf{1 sg.} & \textbf{2 sg.} & \textbf{3 sg.} & \textbf{All pl.} \\
\midrule
lǣdan & lǣde & lǣtst & lǣtt & lǣdað \\
īecan & īece & īecst & īecð & īecað \\
geręstan & geręste & geręstest & geręst & geręstað \\
mētan & mēte & mētst & mētt & mētað \\
gierwan & gierwe & gierest & giereð & gierwað \\
\bottomrule
\end{tabular}
\par\medskip
% panel:Preterite indicative
\textbf{Preterite indicative}\par\smallskip
\setlength{\LTleft}{0pt}\setlength{\LTright}{\fill}
\noindent
\begin{tabular}{@{}>{\raggedright\arraybackslash}p{28mm}>{\raggedright\arraybackslash}p{28mm}>{\raggedright\arraybackslash}p{33mm}>{\raggedright\arraybackslash}p{26mm}@{}}
\toprule
\textbf{Verb} & \textbf{1/\allowbreak{}3 sg.} & \textbf{2 sg.} & \textbf{All pl.} \\
\midrule
lǣdan & lǣdde & lǣddest & lǣddon \\
īecan & īecte & īectest & īecton \\
geręstan & geręste & geręstest & geręston \\
mētan & mētte & mēttest & mētton \\
gierwan & gierede & gieredest & gieredon \\
\bottomrule
\end{tabular}
\par\medskip
\end{minipage}\par

% mdtable:T03:end
% mdparagraph:P016:begin
Keep these contrasts together:
% mdparagraph:P016:end

\begin{itemize}
% mdparagraph:P017:begin
\item After consonant + \textbf{d}, \textbf{sęndan → sęnde}; after vowel + \textbf{d}, \textbf{lǣdan → lǣdde}.
% mdparagraph:P017:end
% mdparagraph:P018:begin
\item After consonant + \textbf{t}, \textbf{geręstan → geręste}; after vowel + \textbf{t}, \textbf{mētan → mētte}. Vowel length does not itself remove a dental.
% mdparagraph:P018:end
% mdparagraph:P019:begin
\item The \textbf{c/\allowbreak{}p/\allowbreak{}x} family has \textbf{t} in the preterite, here \textbf{īecan → īecte}. The exceptional \textbf{gewyrcean → geworhte} below must be learned separately.
% mdparagraph:P019:end
% mdparagraph:P020:begin
\item Before the present second/\allowbreak{}third singular endings, dental combinations give \textbf{sęntst/\allowbreak{}sęnt}, \textbf{lǣtst/\allowbreak{}lǣtt}, and \textbf{mētst/\allowbreak{}mētt}. Final \textbf{tt} can also be written \textbf{t}: \textbf{sęt, lǣt, mēt} are valid third-singular alternatives.
% mdparagraph:P020:end
% mdparagraph:P021:begin
\item \textbf{Geręstest} is our selected second singular present; contracted \textbf{geręst} is also valid. The selected third singular is \textbf{geręst}. Do not add another \textbf{t} to the preterite \textbf{geręste}.
% mdparagraph:P021:end
% mdparagraph:P022:begin
\item In the \textbf{rw/\allowbreak{}lw} family, \textbf{gierwan → gierede}. Our present row retains \textbf{w} in the first singular and plural, but uses \textbf{gierest, giereð} in the other singular persons. The sources also allow present forms with \textbf{w} retained or omitted; the selected row is not a rule deleting every present \textbf{w}.
% mdparagraph:P022:end
\end{itemize}

\reserve{7\baselineskip}
% mdparagraph:P023:begin
\section*{4. Class I: irregular and short stems}
% mdparagraph:P023:end

% teaching: NEW-U02-GRAM-weak-I-irregular
% mdparagraph:P024:begin
Learn the past stems of \textbf{sęllan} and \textbf{gewyrcean} with their meanings. Neither is predicted by the regular long-stem pattern. The lexical prefix \textbf{ge-} is retained in the forms of \textbf{gewyrcean}; no general meaning or tense rule for that prefix is implied. [5]
% mdparagraph:P024:end

% teaching: NEW-U02-PAR-sellan NEW-U02-PAR-gewyrcean
% mdtable:T04:begin
\par\begin{minipage}{\linewidth}\raggedright
\fontsize{10.5}{13.2}\selectfont
% panel:Present indicative
\textbf{Present indicative}\par\smallskip
\setlength{\LTleft}{0pt}\setlength{\LTright}{\fill}
\noindent
\begin{tabular}{@{}>{\raggedright\arraybackslash}p{25.5mm}>{\raggedright\arraybackslash}p{21mm}>{\raggedright\arraybackslash}p{22mm}>{\raggedright\arraybackslash}p{22mm}>{\raggedright\arraybackslash}p{22mm}@{}}
\toprule
\textbf{Verb} & \textbf{1 sg.} & \textbf{2 sg.} & \textbf{3 sg.} & \textbf{All pl.} \\
\midrule
sęllan & sęlle & sęlst & sęlð & sęllað \\
gewyrcean & gewyrce & gewyrcst & gewyrcð & gewyrceað \\
\bottomrule
\end{tabular}
\par\medskip
% panel:Preterite indicative
\textbf{Preterite indicative}\par\smallskip
\setlength{\LTleft}{0pt}\setlength{\LTright}{\fill}
\noindent
\begin{tabular}{@{}>{\raggedright\arraybackslash}p{28mm}>{\raggedright\arraybackslash}p{28mm}>{\raggedright\arraybackslash}p{33mm}>{\raggedright\arraybackslash}p{26mm}@{}}
\toprule
\textbf{Verb} & \textbf{1/\allowbreak{}3 sg.} & \textbf{2 sg.} & \textbf{All pl.} \\
\midrule
sęllan & sealde & sealdest & sealdon \\
gewyrcean & geworhte & geworhtest & geworhton \\
\bottomrule
\end{tabular}
\par\medskip
\end{minipage}\par

% mdtable:T04:end
% mdparagraph:P025:begin
Cook also prints \textbf{sęlest, sęleð} beside \textbf{sęlst, sęlð}, and \textbf{gewyrcan/\allowbreak{}gewyrcað} beside the infinitive \textbf{gewyrcean} and plural \textbf{gewyrceað}. These alternatives do not change the past \textbf{geworhte}.
% mdparagraph:P025:end

% teaching: NEW-U02-GRAM-weak-I-short NEW-U02-GRAM-weak-I-ian
% mdparagraph:P026:begin
The short-stem families have a linking \textbf{e} in the preterite. \textbf{Gefręmman} retains \textbf{mm} in the present first singular and plural, but has single \textbf{m} before \textbf{-est, -eð, -ede}. The consonant + sonorant type \textbf{ęfnan} has \textbf{ęfnede}. Short Class I \textbf{nęrian} retains \textbf{i} in the present first singular and plural, but has \textbf{nęrest, nęreð, nęrede}.
% mdparagraph:P026:end

% teaching: NEW-U02-PAR-gefremman NEW-U02-PAR-efnan NEW-U02-PAR-nerian
% mdtable:T05:begin
\par\begin{minipage}{\linewidth}\raggedright
\fontsize{10.5}{13.2}\selectfont
% panel:Present indicative
\textbf{Present indicative}\par\smallskip
\setlength{\LTleft}{0pt}\setlength{\LTright}{\fill}
\noindent
\begin{tabular}{@{}>{\raggedright\arraybackslash}p{25.5mm}>{\raggedright\arraybackslash}p{21mm}>{\raggedright\arraybackslash}p{22mm}>{\raggedright\arraybackslash}p{22mm}>{\raggedright\arraybackslash}p{22mm}@{}}
\toprule
\textbf{Verb} & \textbf{1 sg.} & \textbf{2 sg.} & \textbf{3 sg.} & \textbf{All pl.} \\
\midrule
gefręmman & gefręmme & gefręmest & gefręmeð & gefręmmað \\
ęfnan & ęfne & ęfnest & ęfneð & ęfnað \\
nęrian & nęrie & nęrest & nęreð & nęriað \\
\bottomrule
\end{tabular}
\par\medskip
% panel:Preterite indicative
\textbf{Preterite indicative}\par\smallskip
\setlength{\LTleft}{0pt}\setlength{\LTright}{\fill}
\noindent
\begin{tabular}{@{}>{\raggedright\arraybackslash}p{28mm}>{\raggedright\arraybackslash}p{28mm}>{\raggedright\arraybackslash}p{33mm}>{\raggedright\arraybackslash}p{26mm}@{}}
\toprule
\textbf{Verb} & \textbf{1/\allowbreak{}3 sg.} & \textbf{2 sg.} & \textbf{All pl.} \\
\midrule
gefręmman & gefręmede & gefręmedest & gefręmedon \\
ęfnan & ęfnede & ęfnedest & ęfnedon \\
nęrian & nęrede & nęredest & nęredon \\
\bottomrule
\end{tabular}
\par\medskip
\end{minipage}\par

% mdtable:T05:end
% mdparagraph:P027:begin
For \textbf{ęfnan}, the contracted past \textbf{ęfnde, ęfndest, ęfndon} is also valid. A source exception for later reference is \textbf{lęcgan}, with past \textbf{lęgde/\allowbreak{}lēde}; it is not a recall item here. Cook’s related example \textbf{hyngran} is impersonal: it must not be substituted into an ordinary personal sentence merely because its endings look familiar.
% mdparagraph:P027:end

\reserve{7\baselineskip}
% mdparagraph:P028:begin
\section*{5. Class II}
% mdparagraph:P028:end

% teaching: NEW-U02-GRAM-weak-II
% mdparagraph:P029:begin
Class II has present \textbf{-ie, -ast, -að, -iað} and the selected past \textbf{-ode, -odest, -odon}. Contrast \textbf{lufast/\allowbreak{}lufode} with Class I \textbf{nęrest/\allowbreak{}nęrede}. Both infinitives end in \textbf{-ian}, but their classes differ. [6]
% mdparagraph:P029:end

% teaching: NEW-U02-PAR-lufian NEW-U02-PAR-gebletsian NEW-U02-PAR-gehalgian
% mdtable:T06:begin
\par\begin{minipage}{\linewidth}\raggedright
\fontsize{10.5}{13.2}\selectfont
% panel:Present indicative
\textbf{Present indicative}\par\smallskip
\setlength{\LTleft}{0pt}\setlength{\LTright}{\fill}
\noindent
\begin{tabular}{@{}>{\raggedright\arraybackslash}p{25.5mm}>{\raggedright\arraybackslash}p{21mm}>{\raggedright\arraybackslash}p{22mm}>{\raggedright\arraybackslash}p{22mm}>{\raggedright\arraybackslash}p{22mm}@{}}
\toprule
\textbf{Verb} & \textbf{1 sg.} & \textbf{2 sg.} & \textbf{3 sg.} & \textbf{All pl.} \\
\midrule
lufian & lufie & lufast & lufað & lufiað \\
geblētsian & geblētsie & geblētsast & geblētsað & geblētsiað \\
gehālgian & gehālgie & gehālgast & gehālgað & gehālgiað \\
\bottomrule
\end{tabular}
\par\medskip
% panel:Preterite indicative
\textbf{Preterite indicative}\par\smallskip
\setlength{\LTleft}{0pt}\setlength{\LTright}{\fill}
\noindent
\begin{tabular}{@{}>{\raggedright\arraybackslash}p{28mm}>{\raggedright\arraybackslash}p{28mm}>{\raggedright\arraybackslash}p{33mm}>{\raggedright\arraybackslash}p{26mm}@{}}
\toprule
\textbf{Verb} & \textbf{1/\allowbreak{}3 sg.} & \textbf{2 sg.} & \textbf{All pl.} \\
\midrule
lufian & lufode & lufodest & lufodon \\
geblētsian & geblētsode & geblētsodest & geblētsodon \\
gehālgian & gehālgode & gehālgodest & gehālgodon \\
\bottomrule
\end{tabular}
\par\medskip
\end{minipage}\par

% mdtable:T06:end
% mdparagraph:P030:begin
Cook also prints \textbf{lufedon}. His class-wide alternatives include past \textbf{-ade, -ude, -ede} and infinitive \textbf{-ig(e)an}; these are reference information, not instructions to generate every spelling for every verb. The dictionary has \textbf{geblēdsian} beside \textbf{geblētsian}. Learn the selected forms above.
% mdparagraph:P030:end

% example: NEW-U02-EXAMPLE-08
% mdparagraph:P031:begin
\textbf{Hē gehālgode ðone dæg.} He hallowed the day.
% mdparagraph:P031:end

\reserve{7\baselineskip}
% mdparagraph:P032:begin
\section*{6. Class III and the negative companion}
% mdparagraph:P032:end

% mdparagraph:P033:begin
All four surviving Class III verbs are included. Their present stem changes and past stems form separate systems; learn each complete row. \textbf{Habban} means possession here, \textbf{libban} living, \textbf{sęcgan} saying, and \textbf{hycgan} thinking. [7]
% mdparagraph:P033:end

% teaching: NEW-U02-GRAM-weak-III-habban NEW-U02-GRAM-weak-III-libban NEW-U02-GRAM-weak-III-secgan NEW-U02-GRAM-weak-III-hycgan
% teaching: NEW-U02-PAR-habban NEW-U02-PAR-libban NEW-U02-PAR-secgan NEW-U02-PAR-hycgan
% mdtable:T07:begin
\par\begin{minipage}{\linewidth}\raggedright
\fontsize{10.5}{13.2}\selectfont
% panel:Present indicative
\textbf{Present indicative}\par\smallskip
\setlength{\LTleft}{0pt}\setlength{\LTright}{\fill}
\noindent
\begin{tabular}{@{}>{\raggedright\arraybackslash}p{25.5mm}>{\raggedright\arraybackslash}p{21mm}>{\raggedright\arraybackslash}p{22mm}>{\raggedright\arraybackslash}p{22mm}>{\raggedright\arraybackslash}p{22mm}@{}}
\toprule
\textbf{Verb} & \textbf{1 sg.} & \textbf{2 sg.} & \textbf{3 sg.} & \textbf{All pl.} \\
\midrule
habban & hæbbe & hæfst & hæfð & habbað \\
libban & libbe & leofast & leofað & libbað \\
sęcgan & sęcge & sægst & sægð & sęcgað \\
hycgan & hycge & hygst & hygð & hycgað \\
\bottomrule
\end{tabular}
\par\medskip
% panel:Preterite indicative
\textbf{Preterite indicative}\par\smallskip
\setlength{\LTleft}{0pt}\setlength{\LTright}{\fill}
\noindent
\begin{tabular}{@{}>{\raggedright\arraybackslash}p{28mm}>{\raggedright\arraybackslash}p{28mm}>{\raggedright\arraybackslash}p{33mm}>{\raggedright\arraybackslash}p{26mm}@{}}
\toprule
\textbf{Verb} & \textbf{1/\allowbreak{}3 sg.} & \textbf{2 sg.} & \textbf{All pl.} \\
\midrule
habban & hæfde & hæfdest & hæfdon \\
libban & lifde & lifdest & lifdon \\
sęcgan & sægde & sægdest & sægdon \\
hycgan & hogode & hogodest & hogodon \\
\bottomrule
\end{tabular}
\par\medskip
\end{minipage}\par

% mdtable:T07:end
% mdparagraph:P034:begin
Reference alternatives include \textbf{hafast, hafað, hæbbað} for \textbf{habban}; \textbf{lifian/\allowbreak{}lifiað} beside \textbf{libban/\allowbreak{}libbað}; \textbf{sęcgean/\allowbreak{}sęcgeað}, \textbf{sęgst/\allowbreak{}sagast}, \textbf{sęgð/\allowbreak{}sagað}, and past \textbf{sǣde, sǣdest, sǣdon} beside the selected \textbf{sęcgan} forms. For \textbf{hycgan}, alternatives include \textbf{hycgean/\allowbreak{}hycgeað}, \textbf{hogast/\allowbreak{}hogað}, and past \textbf{hogde, hogdest, hogdon}. They are valid forms, not extra required paradigms.
% mdparagraph:P034:end

% teaching: NEW-U02-GRAM-nabban-reference
% mdparagraph:P035:begin
\textbf{Nabban}, “have not”, is the contracted negative companion of \textbf{habban}. Recognize its complete system; independent production is not required.
% mdparagraph:P035:end

% teaching: NEW-U02-PAR-nabban
% mdtable:T08:begin
\par\begin{minipage}{\linewidth}\raggedright
\fontsize{10.5}{13.2}\selectfont
% panel:Present indicative
\textbf{Present indicative}\par\smallskip
\setlength{\LTleft}{0pt}\setlength{\LTright}{\fill}
\noindent
\begin{tabular}{@{}>{\raggedright\arraybackslash}p{25.5mm}>{\raggedright\arraybackslash}p{21mm}>{\raggedright\arraybackslash}p{22mm}>{\raggedright\arraybackslash}p{22mm}>{\raggedright\arraybackslash}p{22mm}@{}}
\toprule
\textbf{Verb} & \textbf{1 sg.} & \textbf{2 sg.} & \textbf{3 sg.} & \textbf{All pl.} \\
\midrule
nabban & næbbe & næfst & næfð & nabbað \\
\bottomrule
\end{tabular}
\par\medskip
% panel:Preterite indicative
\textbf{Preterite indicative}\par\smallskip
\setlength{\LTleft}{0pt}\setlength{\LTright}{\fill}
\noindent
\begin{tabular}{@{}>{\raggedright\arraybackslash}p{28mm}>{\raggedright\arraybackslash}p{28mm}>{\raggedright\arraybackslash}p{33mm}>{\raggedright\arraybackslash}p{26mm}@{}}
\toprule
\textbf{Verb} & \textbf{1/\allowbreak{}3 sg.} & \textbf{2 sg.} & \textbf{All pl.} \\
\midrule
nabban & næfde & næfdest & næfdon \\
\bottomrule
\end{tabular}
\par\medskip
\end{minipage}\par

% mdtable:T08:end
% mdparagraph:P036:begin
\textbf{Nafast, nafað} are further source alternatives. The contraction after \textbf{ne} is frequent, not compulsory: uncontracted \textbf{ne hæbbe} is also valid. Recognition of this series does not teach negative forms of the copula.
% mdparagraph:P036:end

\reserve{7\baselineskip}
% mdparagraph:P037:begin
\section*{7. Using the forms}
% mdparagraph:P037:end

% mdparagraph:P038:begin
The following examples are newly composed teaching material.
% mdparagraph:P038:end

% teaching: NEW-U02-GRAM-active-agreement NEW-U02-GRAM-direct-object
% mdparagraph:P039:begin
The finite verb agrees in person and number with its nominative subject. An accusative object does not control the verb. Use the object constructions specified in the vocabulary; a weak paradigm alone does not tell us which case every verb governs. [8]
% mdparagraph:P039:end

% example: NEW-U02-EXAMPLE-01
% mdparagraph:P040:begin
\textbf{Wē hīerað ðone mann.} We hear the man.
% mdparagraph:P040:end

% example: NEW-U02-EXAMPLE-02
% mdparagraph:P041:begin
\textbf{Se mann hīerð ūs.} The man hears us.
% mdparagraph:P041:end

% mdparagraph:P042:begin
In the first example, \textbf{wē} requires a plural verb. In the second, singular \textbf{se mann} requires \textbf{hīerð}, even though its object \textbf{ūs} is plural. Subject–verb–object order is a useful model here; it is not the only possible Old English order.
% mdparagraph:P042:end

% teaching: NEW-U02-GRAM-giving-recipient
% mdparagraph:P043:begin
With \textbf{sęllan}, the thing given is accusative and the recipient is dative. A personal pronoun takes the same case as the noun it replaces.
% mdparagraph:P043:end

% example: NEW-U02-EXAMPLE-03
% mdparagraph:P044:begin
\textbf{Ðū sęlst ðǣm fugle sǣd.} You give the bird seed.
% mdparagraph:P044:end

% teaching: NEW-U02-GRAM-libban-location
% mdparagraph:P045:begin
\textbf{Libban} is intransitive. A location can be expressed by the already learned \textbf{on + dative}. Possessive \textbf{his, hiere, hiera} remain unchanged before the noun.
% mdparagraph:P045:end

% example: NEW-U02-EXAMPLE-04
% mdparagraph:P046:begin
\textbf{Hiera fuglas libbað on his hofe.} Their birds live in his dwelling.
% mdparagraph:P046:end

% teaching: NEW-U02-GRAM-reflexive-use
% mdparagraph:P047:begin
Ordinary personal pronouns also serve reflexively, referring to the subject. With \textbf{geręstan}, use an accusative referring to the person resting. Natural English often omits “myself”, “himself”, and so on in translating this construction. With \textbf{sęllan}, a self-recipient is dative, just like any other recipient. Reflexive meaning comes from reference and construction; there is no special reflexive form in the table. [9]
% mdparagraph:P047:end

% example: NEW-U02-EXAMPLE-05
% mdparagraph:P048:begin
\textbf{Hēo geręste hīe.} She rested herself.
% mdparagraph:P048:end

% teaching: NEW-U02-VOC-ne NEW-U02-GRAM-simple-negation
% mdparagraph:P049:begin
\textbf{Ne} means “not”. In these simple clauses put it immediately before the finite verb. English may need “do not” or “did not”; Old English does not add a corresponding auxiliary. The exercises use noncontracting verbs for negative composition. [10]
% mdparagraph:P049:end

% example: NEW-U02-EXAMPLE-06
% mdparagraph:P050:begin
\textbf{Git ne hīerað ðone mann.} You two do not hear the man.
% mdparagraph:P050:end

% teaching: NEW-U02-GRAM-clause-coordination NEW-U02-GRAM-direct-quotation
% mdparagraph:P051:begin
\textbf{And} joins clauses or predicates. When two coordinated verbs have the same subject, that subject may remain understood with the second verb. Keep each verb’s object and tense clear.
% mdparagraph:P051:end

% mdparagraph:P052:begin
\textbf{Sęcgan} can introduce direct speech: the words inside quotation marks form their own independent utterance. Pronouns inside the quotation refer to that utterance’s speaker and addressees. The construction below does not introduce indirect speech. [11]
% mdparagraph:P052:end

% example: NEW-U02-EXAMPLE-07
% mdparagraph:P053:begin
\textbf{Hē sægð: “Wē habbað wæter.”} He says: “We have water.”
% mdparagraph:P053:end

\reserve{7\baselineskip}
% mdparagraph:P054:begin
\section*{8. Vocabulary and principal parts}
% mdparagraph:P054:end

% mdparagraph:P055:begin
The nouns and the article are those of Unit 1. The following verb senses and constructions are used here. The four entries marked \textbf{forms only} are practised as morphology; their complement constructions are not assigned.
% mdparagraph:P055:end

% teaching: NEW-U02-VOC-hieran NEW-U02-VOC-gefyllan NEW-U02-VOC-cyssan NEW-U02-VOC-settan NEW-U02-VOC-sendan NEW-U02-VOC-laedan NEW-U02-VOC-iecan NEW-U02-VOC-gerestan NEW-U02-VOC-metan NEW-U02-VOC-gierwan NEW-U02-VOC-sellan NEW-U02-VOC-gewyrcean NEW-U02-VOC-gefremman NEW-U02-VOC-efnan NEW-U02-VOC-nerian NEW-U02-VOC-lufian NEW-U02-VOC-gebletsian NEW-U02-VOC-gehalgian NEW-U02-VOC-habban NEW-U02-VOC-libban NEW-U02-VOC-secgan NEW-U02-VOC-hycgan NEW-U02-VOC-nabban
% mdtable:T09:begin
\par\begingroup
\fontsize{10.5}{13.2}\selectfont
\setlength{\LTleft}{0pt}\setlength{\LTright}{\fill}
\begin{longtable}{@{}>{\raggedright\arraybackslash}p{25mm}>{\raggedright\arraybackslash}p{11mm}>{\raggedright\arraybackslash}p{81mm}@{}}
\toprule
\textbf{Verb} & \textbf{Class} & \textbf{Meaning and use} \\
\midrule
\endfirsthead
\toprule
\textbf{Verb} & \textbf{Class} & \textbf{Meaning and use} \\
\midrule
\endhead
\midrule
\multicolumn{3}{r}{\footnotesize Continued on the following page.}\\
\endfoot
\bottomrule
\endlastfoot
hīeran & I & hear; accusative person or thing heard \\
gefyllan & I & finish, accomplish a work; accusative work; no substance-filling construction \\
cyssan & I & kiss; accusative person \\
sęttan & I & set, put down, place; accusative thing; no new direction phrase \\
sęndan & I & send; accusative person or thing; recipient and destination constructions deferred \\
lǣdan & I & lead; accusative person; destination constructions deferred \\
īecan & I & increase; forms only \\
geręstan & I & rest, repose; accusative self-referential pronoun \\
mētan & I & meet, find; accusative person or thing; no “find out” clause \\
gierwan & I & prepare; forms only \\
sęllan & I & give; accusative thing and dative recipient \\
gewyrcean & I & make, build; accusative thing \\
gefręmman & I & perform, accomplish a work; accusative work \\
ęfnan & I & perform; forms only \\
nęrian & I & save; accusative person or thing \\
lufian & II & love; accusative person or thing \\
geblētsian & II & bless; accusative person or thing \\
gehālgian & II & hallow, consecrate, make holy; accusative person or thing \\
habban & III & have, possess; accusative thing possessed; no auxiliary perfect \\
libban & III & live; intransitive, with on + dative for location \\
sęcgan & III & say; direct quotation of an independent utterance \\
hycgan & III & think; forms only \\
nabban & III & have not; complete finite series for recognition only \\
\end{longtable}
\par\endgroup

% mdtable:T09:end
% mdparagraph:P056:begin
The prefixed verbs are separate vocabulary entries. Do not infer a new meaning by freely adding or removing \textbf{ge-}. In particular, \textbf{gefręmman} here means “perform/\allowbreak{}accomplish”, and \textbf{gefyllan} “finish/\allowbreak{}accomplish a work”. The headword variant \textbf{hȳran} belongs with \textbf{hīeran}. [12]
% mdparagraph:P056:end

% teaching: NEW-U02-GRAM-citation-boundary
% mdparagraph:P057:begin
The three principal parts are supplied below for lookup. The middle column is the learned first/\allowbreak{}third singular preterite indicative. The infinitive and participle columns identify dictionary forms; they are not independent recall targets and do not teach infinitival, passive, perfect or participial syntax. Parentheses reproduce optional material in the source citation, not a freely applicable prefix rule.
% mdparagraph:P057:end

% mdtable:T10:begin
\par\begingroup
\fontsize{10.5}{13.2}\selectfont
\setlength{\LTleft}{0pt}\setlength{\LTright}{\fill}
\begin{longtable}{@{}>{\raggedright\arraybackslash}p{37mm}>{\raggedright\arraybackslash}p{36mm}>{\raggedright\arraybackslash}p{44mm}@{}}
\toprule
\textbf{Infinitive citation} & \textbf{Preterite 1/\allowbreak{}3 sg.} & \textbf{Past participle citation} \\
\midrule
\endfirsthead
\toprule
\textbf{Infinitive citation} & \textbf{Preterite 1/\allowbreak{}3 sg.} & \textbf{Past participle citation} \\
\midrule
\endhead
\midrule
\multicolumn{3}{r}{\footnotesize Continued on the following page.}\\
\endfoot
\bottomrule
\endlastfoot
hīeran & hīerde & (ge)hīered \\
gefyllan & gefylde & gefylled \\
cyssan & cyste & (ge)cyssed \\
sęttan & sętte & (ge)sętt \\
sęndan & sęnde & (ge)sęnd(ed) \\
lǣdan & lǣdde & (ge)lǣd(ed) \\
īecan & īecte & (ge)īeced \\
geręstan & geręste & geręst \\
mētan & mētte & (ge)mēt(t) \\
gierwan & gierede & (ge)gier(w)ed \\
sęllan & sealde & (ge)seald \\
gewyrcean & geworhte & geworht \\
gefręmman & gefręmede & gefręmed \\
ęfnan & ęfnede & (ge)ęfned \\
nęrian & nęrede & (ge)nęred \\
lufian & lufode & (ge)lufod \\
geblētsian & geblētsode & geblētsod \\
gehālgian & gehālgode & gehālgod \\
habban & hæfde & (ge)hæfd \\
libban & lifde & (ge)lifd \\
sęcgan & sægde & (ge)sægd \\
hycgan & hogode & (ge)hogod \\
nabban & næfde & (ge)næfd \\
\end{longtable}
\par\endgroup

% mdtable:T10:end
% mdparagraph:P058:begin
The alternative past series of \textbf{sęcgan} and \textbf{hycgan} have the participle citations \textbf{(ge)sǣd} and \textbf{(ge)hogd}. These too are reference forms only.
% mdparagraph:P058:end

\reserve{7\baselineskip}
% mdparagraph:P059:begin
\section*{9. Pronoun practice}
% mdparagraph:P059:end

% mdparagraph:P060:begin
Use the selected forms in §1. Identical spellings must still be labelled with their separate cases.
% mdparagraph:P060:end

\par\begin{minipage}{\linewidth}\raggedright
% item: NEW-U02-P01
% mdparagraph:P061:begin
\itempara{1}{Write the nominative, dative and accusative personal pronouns for the first person singular and the second person singular. Give the English meaning of each row.}
% mdparagraph:P061:end

\end{minipage}\par\addvspace{4pt}

\par\begin{minipage}{\linewidth}\raggedright
% item: NEW-U02-P02
% mdparagraph:P062:begin
\itempara{2}{Write the nominative, dative and accusative personal pronouns for “we two” and “you two”. State which number of verb form goes with each nominative.}
% mdparagraph:P062:end

\end{minipage}\par\addvspace{4pt}

\par\begin{minipage}{\linewidth}\raggedright
% item: NEW-U02-P03
% mdparagraph:P063:begin
\itempara{3}{Write the nominative, dative and accusative personal pronouns for “we” and “you” when more than two people are meant.}
% mdparagraph:P063:end

\end{minipage}\par\addvspace{4pt}

\par\begin{minipage}{\linewidth}\raggedright
% item: NEW-U02-P04
% mdparagraph:P064:begin
\itempara{4}{Write the third-person singular nominative, dative and accusative pronouns for masculine, feminine and neuter referents. Keep the three gender rows distinct.}
% mdparagraph:P064:end

\end{minipage}\par\addvspace{4pt}

\par\begin{minipage}{\linewidth}\raggedright
% item: NEW-U02-P05
% mdparagraph:P065:begin
\itempara{5}{Write the third-person plural nominative, dative and accusative personal pronouns. Are separate forms needed for the three genders?}
% mdparagraph:P065:end

\end{minipage}\par\addvspace{4pt}

\par\begin{minipage}{\linewidth}\raggedright
% item: NEW-U02-P06
% mdparagraph:P066:begin
\itempara{6}{Match these genitive forms to person and number: mīn, ðīn, uncer, incer, ūre, ēower. Use “first” or “second” person and “singular”, “dual” or “plural”.}
% mdparagraph:P066:end

\end{minipage}\par\addvspace{4pt}

\reserve{7\baselineskip}
% mdparagraph:P067:begin
\section*{10. Verb forms}
% mdparagraph:P067:end

% mdparagraph:P068:begin
Except for the given negative forms in exercise 19 and the lookup exercise 20, produce the requested forms from memory. When a complete paradigm is requested, label all three singular and all three plural persons in both tenses; identical plural forms may be grouped.
% mdparagraph:P068:end

\par\begin{minipage}{\linewidth}\raggedright
% item: NEW-U02-F01
% mdparagraph:P069:begin
\itempara{1}{Write the complete present and preterite indicative of hīeran. Group identical plural forms if you label all three persons. Then give “you two heard”.}
% mdparagraph:P069:end

\end{minipage}\par\addvspace{4pt}

\par\begin{minipage}{\linewidth}\raggedright
% item: NEW-U02-F02
% mdparagraph:P070:begin
\itempara{2}{Give “you (singular) finish”, “he finishes”, “I finished” and “we finished”, using gefyllan in the sense “finish; accomplish”. Explain the difference between the present and preterite stem spellings.}
% mdparagraph:P070:end

\end{minipage}\par\addvspace{4pt}

\par\begin{minipage}{\linewidth}\raggedright
% item: NEW-U02-F03
% mdparagraph:P071:begin
\itempara{3}{Give the second and third singular present indicative and the second singular and plural preterite indicative of cyssan.}
% mdparagraph:P071:end

\end{minipage}\par\addvspace{4pt}

\par\begin{minipage}{\linewidth}\raggedright
% item: NEW-U02-F04
% mdparagraph:P072:begin
\itempara{4}{For sęttan and sęndan, give the first singular present, first singular preterite and plural preterite. Label the tense of each form.}
% mdparagraph:P072:end

\end{minipage}\par\addvspace{4pt}

\par\begin{minipage}{\linewidth}\raggedright
% item: NEW-U02-F05
% mdparagraph:P073:begin
\itempara{5}{For lǣdan and mētan, give the second singular present, third singular present and second singular preterite. Keep every printed length mark.}
% mdparagraph:P073:end

\end{minipage}\par\addvspace{4pt}

\par\begin{minipage}{\linewidth}\raggedright
% item: NEW-U02-F06
% mdparagraph:P074:begin
\itempara{6}{For īecan and gierwan, give the first singular present, third singular present and first singular preterite. This is a form exercise, not a sentence translation.}
% mdparagraph:P074:end

\end{minipage}\par\addvspace{4pt}

\par\begin{minipage}{\linewidth}\raggedright
% item: NEW-U02-F07
% mdparagraph:P075:begin
\itempara{7}{Write the complete present and preterite indicative of geręstan. Then give “he rested himself” using its ordinary reflexive construction.}
% mdparagraph:P075:end

\end{minipage}\par\addvspace{4pt}

\par\begin{minipage}{\linewidth}\raggedright
% item: NEW-U02-F08
% mdparagraph:P076:begin
\itempara{8}{Write the complete present and preterite indicative of sęllan. Then give “I give” and “I gave” with the personal pronoun.}
% mdparagraph:P076:end

\end{minipage}\par\addvspace{4pt}

\par\begin{minipage}{\linewidth}\raggedright
% item: NEW-U02-F09
% mdparagraph:P077:begin
\itempara{9}{Give the second and third singular present indicative and all persons of the singular preterite indicative of gewyrcean.}
% mdparagraph:P077:end

\end{minipage}\par\addvspace{4pt}

\par\begin{minipage}{\linewidth}\raggedright
% item: NEW-U02-F10
% mdparagraph:P078:begin
\itempara{10}{For gefręmman, give the first, second and third singular present, the plural present, and the first singular preterite.}
% mdparagraph:P078:end

\end{minipage}\par\addvspace{4pt}

\par\begin{minipage}{\linewidth}\raggedright
% item: NEW-U02-F11
% mdparagraph:P079:begin
\itempara{11}{For ęfnan, give the first singular present, second singular present, first singular preterite and plural preterite.}
% mdparagraph:P079:end

\end{minipage}\par\addvspace{4pt}

\par\begin{minipage}{\linewidth}\raggedright
% item: NEW-U02-F12
% mdparagraph:P080:begin
\itempara{12}{Write the complete present and preterite indicative of nęrian. Then give “we two save” and “you (singular) saved”.}
% mdparagraph:P080:end

\end{minipage}\par\addvspace{4pt}

\par\begin{minipage}{\linewidth}\raggedright
% item: NEW-U02-F13
% mdparagraph:P081:begin
\itempara{13}{Write the complete present and preterite indicative of lufian. Identify the vowel in the present second singular ending and in the preterite stem.}
% mdparagraph:P081:end

\end{minipage}\par\addvspace{4pt}

\par\begin{minipage}{\linewidth}\raggedright
% item: NEW-U02-F14
% mdparagraph:P082:begin
\itempara{14}{For geblētsian and gehālgian, give the third singular present, second singular preterite and plural preterite.}
% mdparagraph:P082:end

\end{minipage}\par\addvspace{4pt}

\par\begin{minipage}{\linewidth}\raggedright
% item: NEW-U02-F15
% mdparagraph:P083:begin
\itempara{15}{Write the complete present and preterite indicative of habban. Include the personal pronouns in the three singular rows.}
% mdparagraph:P083:end

\end{minipage}\par\addvspace{4pt}

\par\begin{minipage}{\linewidth}\raggedright
% item: NEW-U02-F16
% mdparagraph:P084:begin
\itempara{16}{Write the complete present and preterite indicative of libban.}
% mdparagraph:P084:end

\end{minipage}\par\addvspace{4pt}

\par\begin{minipage}{\linewidth}\raggedright
% item: NEW-U02-F17
% mdparagraph:P085:begin
\itempara{17}{Write the complete present and preterite indicative of sęcgan.}
% mdparagraph:P085:end

\end{minipage}\par\addvspace{4pt}

\par\begin{minipage}{\linewidth}\raggedright
% item: NEW-U02-F18
% mdparagraph:P086:begin
\itempara{18}{Write the complete present and preterite indicative of hycgan.}
% mdparagraph:P086:end

\end{minipage}\par\addvspace{4pt}

\par\begin{minipage}{\linewidth}\raggedright
% item: NEW-U02-F19
% mdparagraph:P087:begin
\itempara{19}{For each given form, identify person, number and tense, allowing more than one answer where necessary: næbbe, næfst, næfð, nabbað, næfde, næfdest, næfdon. Give the negative meaning.}
% mdparagraph:P087:end

\end{minipage}\par\addvspace{4pt}

\par\begin{minipage}{\linewidth}\raggedright
% item: NEW-U02-F20
% mdparagraph:P088:begin
\itempara{20}{Use the principal-part entries to identify the lemma and tense of hīerde, sealde, geworhte, gefręmede, lufode and lifde. Which of these forms can be first or third singular in the indicative?}
% mdparagraph:P088:end

\end{minipage}\par\addvspace{4pt}

\reserve{7\baselineskip}
% mdparagraph:P089:begin
\section*{11. Cases, reference and clauses}
% mdparagraph:P089:end

% mdparagraph:P090:begin
Use the meanings and constructions in §8. Give complete sentences where a translation is requested.
% mdparagraph:P090:end

\par\begin{minipage}{\linewidth}\raggedright
% item: NEW-U02-C01
% mdparagraph:P091:begin
\itempara{1}{The referents are a man and several birds. Translate Hē hīerð hīe and Hīe hīerað hine. Identify the case and referent number of every pronoun, and explain how the two statements differ.}
% mdparagraph:P091:end

\end{minipage}\par\addvspace{4pt}

\par\begin{minipage}{\linewidth}\raggedright
% item: NEW-U02-C02
% mdparagraph:P092:begin
\itempara{2}{Translate “I give the man seed”. Then replace only “the man” with the appropriate personal pronoun, keeping the same meaning.}
% mdparagraph:P092:end

\end{minipage}\par\addvspace{4pt}

\par\begin{minipage}{\linewidth}\raggedright
% item: NEW-U02-C03
% mdparagraph:P093:begin
\itempara{3}{Translate Wit habbað męte. Se mann sęlð unc wæter. Identify the case and number of wit and unc.}
% mdparagraph:P093:end

\end{minipage}\par\addvspace{4pt}

\par\begin{minipage}{\linewidth}\raggedright
% item: NEW-U02-C04
% mdparagraph:P094:begin
\itempara{4}{Translate “You two save me. I hear you two.” Use the dual personal pronouns where appropriate.}
% mdparagraph:P094:end

\end{minipage}\par\addvspace{4pt}

\par\begin{minipage}{\linewidth}\raggedright
% item: NEW-U02-C05
% mdparagraph:P095:begin
\itempara{5}{Replace sēo eorðe in the second sentence with the appropriate personal pronoun, and translate both sentences: Sēo eorðe hæfð wæter. Sēo eorðe hæfð gærs.}
% mdparagraph:P095:end

\end{minipage}\par\addvspace{4pt}

\par\begin{minipage}{\linewidth}\raggedright
% item: NEW-U02-C06
% mdparagraph:P096:begin
\itempara{6}{Translate “His dwelling is on the earth. His birds live in the dwelling.” The possessor is a man in both sentences.}
% mdparagraph:P096:end

\end{minipage}\par\addvspace{4pt}

\par\begin{minipage}{\linewidth}\raggedright
% item: NEW-U02-C07
% mdparagraph:P097:begin
\itempara{7}{Translate “She has a dwelling. Birds live in her dwelling.”}
% mdparagraph:P097:end

\end{minipage}\par\addvspace{4pt}

\par\begin{minipage}{\linewidth}\raggedright
% item: NEW-U02-C08
% mdparagraph:P098:begin
\itempara{8}{Translate “They have a dwelling. Their dwelling has birds.”}
% mdparagraph:P098:end

\end{minipage}\par\addvspace{4pt}

\par\begin{minipage}{\linewidth}\raggedright
% item: NEW-U02-C09
% mdparagraph:P099:begin
\itempara{9}{Translate “I rested myself. He rested himself.” Mark the subject and object in each clause.}
% mdparagraph:P099:end

\end{minipage}\par\addvspace{4pt}

\par\begin{minipage}{\linewidth}\raggedright
% item: NEW-U02-C10
% mdparagraph:P100:begin
\itempara{10}{Translate “You two rest yourselves. They rest themselves.”}
% mdparagraph:P100:end

\end{minipage}\par\addvspace{4pt}

\par\begin{minipage}{\linewidth}\raggedright
% item: NEW-U02-C11
% mdparagraph:P101:begin
\itempara{11}{Change Ic hīere ðone fugol and Hē lufað ðā fuglas into negative statements. Translate the results.}
% mdparagraph:P101:end

\end{minipage}\par\addvspace{4pt}

\par\begin{minipage}{\linewidth}\raggedright
% item: NEW-U02-C12
% mdparagraph:P102:begin
\itempara{12}{Translate “You (singular) did not hear him. We did not send the birds.”}
% mdparagraph:P102:end

\end{minipage}\par\addvspace{4pt}

\par\begin{minipage}{\linewidth}\raggedright
% item: NEW-U02-C13
% mdparagraph:P103:begin
\itempara{13}{Translate “I give myself food. I give him food.” Explain which pronoun is reflexive and why both are dative.}
% mdparagraph:P103:end

\end{minipage}\par\addvspace{4pt}

\reserve{7\baselineskip}
% mdparagraph:P104:begin
\section*{12. Translation}
% mdparagraph:P104:end

\par\begin{minipage}{\linewidth}\raggedright
% item: NEW-U02-T01
% mdparagraph:P105:begin
\itempara{1}{Translate “The man put down the seed. The men sent it.”}
% mdparagraph:P105:end

\end{minipage}\par\addvspace{4pt}

\par\begin{minipage}{\linewidth}\raggedright
% item: NEW-U02-T02
% mdparagraph:P106:begin
\itempara{2}{Translate “I met you (singular). You met me. We two rested ourselves.”}
% mdparagraph:P106:end

\end{minipage}\par\addvspace{4pt}

\par\begin{minipage}{\linewidth}\raggedright
% item: NEW-U02-T03
% mdparagraph:P107:begin
\itempara{3}{Translate “He saved the man and loved him. He did not send him.”}
% mdparagraph:P107:end

\end{minipage}\par\addvspace{4pt}

\par\begin{minipage}{\linewidth}\raggedright
% item: NEW-U02-T04
% mdparagraph:P108:begin
\itempara{4}{Translate “You (singular) blessed the day. They hallowed it.”}
% mdparagraph:P108:end

\end{minipage}\par\addvspace{4pt}

\par\begin{minipage}{\linewidth}\raggedright
% item: NEW-U02-T05
% mdparagraph:P109:begin
\itempara{5}{Translate Hē geworhte his hof. Hīe gefręmedon hiera weorc. Hē geręste hine on ðǣm hofe.}
% mdparagraph:P109:end

\end{minipage}\par\addvspace{4pt}

\par\begin{minipage}{\linewidth}\raggedright
% item: NEW-U02-T06
% mdparagraph:P110:begin
\itempara{6}{Translate “We have the man’s seed. We give him water. He hears us.”}
% mdparagraph:P110:end

\end{minipage}\par\addvspace{4pt}

\reserve{7\baselineskip}
% mdparagraph:P111:begin
\section*{13. Two connected readings}
% mdparagraph:P111:end

% mdparagraph:P112:begin
Both readings are new compositions using the forms and constructions of Units 1–2.
% mdparagraph:P112:end

\reserve{7\baselineskip}
% mdparagraph:P113:begin
\subsection*{Seed and water}
% mdparagraph:P113:end

% composition: NEW-U02-TEXT-NARRATIVE-01
% mdparagraph:P114:begin
\par\begin{minipage}{\linewidth}\raggedright
\readingtext{Se mann hæfde sǣd and wæter. Ðā męnn hæfdon hof and fuglas. Ðā fuglas lifdon on hiera hofe. Se mann sealde ðǣm mannum sǣd and wæter. Hīe lufodon ðā fuglas and hīerdon hīe. Hīe sealdon him sǣd and wæter. Hīe ne sęndon ðā fuglas. Hīe gefręmedon hiera weorc and geręston hīe. Se mann sægde: “Ic lufie ðā fuglas.”}
\end{minipage}\par
% mdparagraph:P114:end

% mdparagraph:P115:begin
Plural pronoun spellings do not always identify a unique antecedent. Follow the sequence of giving and caring for the birds, and be prepared to explain any other defensible reference.
% mdparagraph:P115:end

\reserve{7\baselineskip}
% mdparagraph:P116:begin
\subsection*{The workers and the birds}
% mdparagraph:P116:end

% mdparagraph:P117:begin
Four workers, A, B, C and D, are together with the birds. Speaker labels identify the addressees. In the last turn, wē includes all four workers and gē addresses B, C and D.
% mdparagraph:P117:end

% composition: NEW-U02-TEXT-DIALOGUE-01
% mdparagraph:P118:begin
\textbf{A to B:} Ic hæbbe sǣd. Ðū hæfst wæter.
% mdparagraph:P118:end

% mdparagraph:P119:begin
\textbf{B to A:} Wit sęllað ðǣm fuglum męte. Hīe hīerað unc.
% mdparagraph:P119:end

% mdparagraph:P120:begin
\textbf{C to A, B:} Git lufiað ðā fuglas. Ic hīere inc.
% mdparagraph:P120:end

% mdparagraph:P121:begin
\textbf{A to B, C, D:} Wē lufiað hīe. Gē hīerað hīe, and hīe hīerað ēow.
% mdparagraph:P121:end

\reserve{7\baselineskip}
% mdparagraph:P122:begin
\subsection*{Reading exercises}
% mdparagraph:P122:end

\par\begin{minipage}{\linewidth}\raggedright
% item: NEW-U02-N01
% mdparagraph:P123:begin
\itempara{1}{Translate the connected account “Seed and water” in full. Keep the sequence of giver, recipients and things given clear.}
% mdparagraph:P123:end

\end{minipage}\par\addvspace{4pt}

\par\begin{minipage}{\linewidth}\raggedright
% item: NEW-U02-N02
% mdparagraph:P124:begin
\itempara{2}{In the account, trace the people and birds referred to by each plural personal pronoun. Which references are settled by the form itself, and which depend on the surrounding actions?}
% mdparagraph:P124:end

\end{minipage}\par\addvspace{4pt}

\par\begin{minipage}{\linewidth}\raggedright
% item: NEW-U02-N03
% mdparagraph:P125:begin
\itempara{3}{Translate the connected conversation in full. In each turn, state who is included in ic, ðū, wit, git, wē and gē.}
% mdparagraph:P125:end

\end{minipage}\par\addvspace{4pt}

\par\begin{minipage}{\linewidth}\raggedright
% item: NEW-U02-N04
% mdparagraph:P126:begin
\itempara{4}{In the conversation, identify the case, number and referents of unc, inc and ēow. In each clause containing one of these objects, identify the subject and explain which pronoun controls the finite verb.}
% mdparagraph:P126:end

\end{minipage}\par\addvspace{4pt}

\par\begin{minipage}{\linewidth}\raggedright
% item: NEW-U02-N05
% mdparagraph:P127:begin
\itempara{5}{Retell the first four sentences of “Seed and water” as present actions or states. Preserve the people, things and cases.}
% mdparagraph:P127:end

\end{minipage}\par\addvspace{4pt}

\par\begin{minipage}{\linewidth}\raggedright
% item: NEW-U02-N06
% mdparagraph:P128:begin
\itempara{6}{In Hīe gefręmedon hiera weorc and geręston hīe, identify the subject, the object of each verb, and the referent of hiera. Translate the whole sentence.}
% mdparagraph:P128:end

\end{minipage}\par\addvspace{4pt}

\reserve{7\baselineskip}
% mdparagraph:P129:begin
\section*{14. Source reading: the Sabbath}
% mdparagraph:P129:end

% mdparagraph:P130:begin
This is the complete closing paragraph of Cook’s “The Creation of the World”, printed p.128 /\allowbreak{} PDF p.144. Its seven prose lines are editorially counted; only the marginal number \textbf{5} is printed. The linguistic text is unchanged; note cues and line breaks are removed. Use the local help freely. The six weak verbs are now learned, but the passive, adjective, temporal, relative, causal and infinitival constructions below are supplied for this passage. [13]
% mdparagraph:P130:end

% source: FB-I-CREATION-SABBATH-01
% mdparagraph:P131:begin
\par\begin{minipage}{\linewidth}\raggedright
\readingtext{Eornostlice ðā wǣron fullfręmede heofonas and eorðe and eall hiera frætwung. And God ðā gefylde on ðone seofoðan dæg his weorc ðe hē geworhte, and hē geręste hine on ðone seofoðan dæg fram eallum ðǣm weorce ðe hē gefręmede. And God geblētsode ðone seofoðan dæg and hine gehālgode, for-ðon-ðe hē on ðone dæg geswāc his weorces ðe hē gescēop tō wyrceanne.}
\end{minipage}\par
% mdparagraph:P131:end

\reserve{7\baselineskip}
% mdparagraph:P132:begin
\subsection*{Local help}
% mdparagraph:P132:end

% teaching: NEW-U02-SOURCE-SUPPORT NEW-U02-GRAM-reading-sabbath-adverbials NEW-U02-GRAM-reading-sabbath-passive
% mdparagraph:P133:begin
\textbf{Eornostlice}, “accordingly, thus”; \textbf{ðā}, “then”, an adverb here rather than the identically spelled article.
% mdparagraph:P133:end

% mdparagraph:P134:begin
\textbf{Wǣron fullfręmede}, “were completed/\allowbreak{}finished”. This is a passive expression. \textbf{Wǣron} is familiar, but its use with the plural participle \textbf{fullfręmede} is supplied. \textbf{Fullfręmede} is not another finite past tense merely because it ends in \textbf{-ede}. The coordinated subject follows the verb: the heavens, the earth, and all their array.
% mdparagraph:P134:end

% teaching: NEW-U02-GRAM-reading-sabbath-adjective-eall NEW-U02-GRAM-reading-sabbath-ordinal NEW-U02-GRAM-reading-sabbath-temporal-on-accusative
% mdparagraph:P135:begin
\textbf{Eall hiera frætwung}, “all their array”: \textbf{eall} modifies feminine nominative singular \textbf{frætwung}. In \textbf{eallum ðǣm weorce}, “all the work”, \textbf{eallum} and \textbf{ðǣm} agree with neuter dative singular \textbf{weorce}. The ending \textbf{-um} does not make this noun phrase plural.
% mdparagraph:P135:end

% mdparagraph:P136:begin
\textbf{Ðone seofoðan dæg}, “the seventh day”: masculine accusative singular, with the supplied weak ordinal form \textbf{seofoðan}. \textbf{On ðone seofoðan dæg} means “on the seventh day”; \textbf{on ðone dæg} means “on that day”, the day just named. These are temporal \textbf{accusative} phrases. They express neither location with the dative nor motion; do not change \textbf{dæg} to \textbf{dæge} by extending the earlier locative rule. [14]
% mdparagraph:P136:end

% teaching: NEW-U02-GRAM-reading-sabbath-relative-object
% mdparagraph:P137:begin
The three relative clauses modify the preceding noun for “work”:
% mdparagraph:P137:end

\begin{itemize}
% mdparagraph:P138:begin
\item \textbf{ðe hē geworhte}, “which he made”, modifies \textbf{weorc}.
% mdparagraph:P138:end
% mdparagraph:P139:begin
\item \textbf{ðe hē gefręmede}, “which he performed”, modifies \textbf{weorce}.
% mdparagraph:P139:end
% mdparagraph:P140:begin
\item \textbf{ðe hē gescēop tō wyrceanne}, “which he created to make”, modifies \textbf{weorces}.
% mdparagraph:P140:end
\end{itemize}

% mdparagraph:P141:begin
In each clause, \textbf{ðe} represents the direct object. The preceding noun has its own case in the larger sentence: accusative, dative and genitive respectively. Relative \textbf{ðe} has no macron; it is not personal \textbf{ðē}. English “had made/\allowbreak{}had performed” may convey the earlier past action in context. General relative syntax is not required. [15]
% mdparagraph:P141:end

% teaching: NEW-U02-GRAM-reading-sabbath-rest-fram
% mdparagraph:P142:begin
\textbf{Geręste hine … fram eallum ðǣm weorce}, “rested … from all the work”. \textbf{Fram} takes dative \textbf{weorce}. This particular rest-from-work construction is supplied; the reflexive pronoun and finite weak verb are learned.
% mdparagraph:P142:end

% teaching: NEW-U02-GRAM-reading-sabbath-causal NEW-U02-GRAM-reading-sabbath-cessation-genitive NEW-U02-GRAM-reading-sabbath-created-to-make
% mdparagraph:P143:begin
\textbf{For-ðon-ðe} introduces the reason, “because”. The causal clause continues to the end of the paragraph.
% mdparagraph:P143:end

% mdparagraph:P144:begin
\textbf{Geswāc his weorces}, “ceased from his work”: \textbf{geswāc} is a supplied strong past form, third singular; the verb here governs genitive \textbf{weorces}. Recognizing the noun’s genitive does not by itself explain this government. [16]
% mdparagraph:P144:end

% mdparagraph:P145:begin
\textbf{Gescēop}, “created”, is another supplied third-singular strong past form. The whole ending \textbf{ðe hē gescēop tō wyrceanne} is “which he created to make”. \textbf{Tō wyrceanne} is an inflected infinitive phrase; it is not a second finite verb. The understood maker is God. Cook explicitly retains this “created to make” relation: preserve it rather than substitute a familiar modern version of the Creation story. A close English rendering may be awkward; no single polished paraphrase is required. [17]
% mdparagraph:P145:end

\reserve{7\baselineskip}
% mdparagraph:P146:begin
\subsection*{Reading exercises}
% mdparagraph:P146:end

\par\begin{minipage}{\linewidth}\raggedright
% item: NEW-U02-R01
% mdparagraph:P147:begin
\itempara{1}{Translate the complete Sabbath passage using its local support. Preserve each action, referent and clause link.}
% mdparagraph:P147:end

\end{minipage}\par\addvspace{4pt}

\par\begin{minipage}{\linewidth}\raggedright
% item: NEW-U02-R02
% mdparagraph:P148:begin
\itempara{2}{Using the vocabulary and principal-part tables for headword lookup, find the six weak preterite indicative forms in the passage. Give their headwords, class, person and number.}
% mdparagraph:P148:end

\end{minipage}\par\addvspace{4pt}

\par\begin{minipage}{\linewidth}\raggedright
% item: NEW-U02-R03
% mdparagraph:P149:begin
\itempara{3}{Compare the two occurrences of hine: geręste hine and hine gehālgode. For each, identify its antecedent and case, translate it naturally, and state whether its use is reflexive.}
% mdparagraph:P149:end

\end{minipage}\par\addvspace{4pt}

\par\begin{minipage}{\linewidth}\raggedright
% item: NEW-U02-R04
% mdparagraph:P150:begin
\itempara{4}{In hiera frætwung, his weorc and his weorces, identify the possessor and the possessed noun. Explain why the pronoun spelling is unchanged when the head’s number or case changes. Use the supplied construction note for weorces.}
% mdparagraph:P150:end

\end{minipage}\par\addvspace{4pt}

\par\begin{minipage}{\linewidth}\raggedright
% item: NEW-U02-R05
% mdparagraph:P151:begin
\itempara{5}{Using the supplied opening construction, identify the case and number of heofonas, eorðe and frætwung. Which complete phrase supplies the subject of wǣron?}
% mdparagraph:P151:end

\end{minipage}\par\addvspace{4pt}

\par\begin{minipage}{\linewidth}\raggedright
% item: NEW-U02-R06
% mdparagraph:P152:begin
\itempara{6}{List in English the actions assigned to God on the seventh day, using the passage’s temporal support.}
% mdparagraph:P152:end

\end{minipage}\par\addvspace{4pt}

\par\begin{minipage}{\linewidth}\raggedright
% item: NEW-U02-R07
% mdparagraph:P153:begin
\itempara{7}{Using the local clause support, explain the reason the passage gives for hallowing the day and how the final creation phrase describes the work.}
% mdparagraph:P153:end

\end{minipage}\par\addvspace{4pt}

\reserve{7\baselineskip}
% mdparagraph:P154:begin
\section*{15. Review}
% mdparagraph:P154:end

\par\begin{minipage}{\linewidth}\raggedright
% item: NEW-U02-X01
% mdparagraph:P155:begin
\itempara{1}{Change this passage to the preterite, then make only its final clause negative: Wit habbað sǣd. Wē sęllað ðǣm fuglum męte. Hīe hīerað ūs.}
% mdparagraph:P155:end

\end{minipage}\par\addvspace{4pt}

\par\begin{minipage}{\linewidth}\raggedright
% item: NEW-U02-X02
% mdparagraph:P156:begin
\itempara{2}{Correct the four forms that prevent this passage from meaning “He gives me seed. We two have it. The birds hear us two”: Hine sęlð ic sǣd. Wit hæfst hit. Ðā fuglas hīerað ūre.}
% mdparagraph:P156:end

\end{minipage}\par\addvspace{4pt}

\par\begin{minipage}{\linewidth}\raggedright
% item: NEW-U02-X03
% mdparagraph:P157:begin
\itempara{3}{Write three connected Old English sentences: say that the man had birds, that he gave them seed, and that he rested himself. Use a personal pronoun after introducing the man.}
% mdparagraph:P157:end

\end{minipage}\par\addvspace{4pt}

\clearpage
% mdparagraph:P158:begin
\section*{Source notes}
% mdparagraph:P158:end

\begingroup\fontsize{10}{12.8}\selectfont
% mdparagraph:P159:begin
[1] Cook, \emph{A First Book in Old English}, second edition (1895), §§81–82, pp.48–49 /\allowbreak{} PDF64–65, supplies the personal forms. Ringe and Taylor, \emph{The Development of Old English} (2014), §1.1.2, p.6 /\allowbreak{} publisher-preview PDF21, explicitly confirms plural verbs with dual subjects. Cook’s spelling is retained.
% mdparagraph:P159:end

% mdparagraph:P160:begin
[2] Cook §83, p.50 /\allowbreak{} PDF66: invariant third-person genitives and the separate inflecting possessive adjectives.
% mdparagraph:P160:end

% mdparagraph:P161:begin
[3] Cook §96, p.56 /\allowbreak{} PDF72, and the complete weak paradigms, §§117–124, pp.72–76 /\allowbreak{} PDF88–92. Some complete rows here combine a checked lexical stem with an explicit source pattern; they do not claim that every displayed cell occurs in the reading.
% mdparagraph:P161:end

% mdparagraph:P162:begin
[4] Cook §113, pp.67–69 /\allowbreak{} PDF83–85, with §§23,33–36, pp.18–19,24–25 /\allowbreak{} PDF34–35,40–41. Independent checks: Wright and Wright, \emph{Old English Grammar}, second edition (1914), §§305,476,524,528–529,533, pp.151,248,267,269–270,273 /\allowbreak{} PDF171,268,287,289–290,293. These corroborate the contractions, inherited consonants and selected \textbf{gierwan} forms; no alternative dialectal spelling is silently imported.
% mdparagraph:P162:end

% mdparagraph:P163:begin
[5] Cook §§114–117, pp.69–72 /\allowbreak{} PDF85–88; §18, p.14 /\allowbreak{} PDF30, for optional \textbf{e} after \textbf{c/\allowbreak{}g} before a further vowel. Wright and Wright §532, p.273 /\allowbreak{} PDF293, confirms the \textbf{ęfnan} type and both past patterns. The short-stem exception and impersonal limit follow Cook §§115,171,190.
% mdparagraph:P163:end

% mdparagraph:P164:begin
[6] Cook §§118–119, pp.73–74 /\allowbreak{} PDF89–90.
% mdparagraph:P164:end

% mdparagraph:P165:begin
[7] Cook §§120–124, pp.74–76 /\allowbreak{} PDF90–92; §29, p.23 /\allowbreak{} PDF39, on frequent negative contraction.
% mdparagraph:P165:end

% mdparagraph:P166:begin
[8] Cook §§168,189–191, pp.96,102 /\allowbreak{} PDF112,118; giving recipients §§160,164(a), pp.93–94 /\allowbreak{} PDF109–110. Lexical constructions remain limited to the senses stated in §8.
% mdparagraph:P166:end

% mdparagraph:P167:begin
[9] Cook §§82,184(a,b), pp.49,100–101 /\allowbreak{} PDF65,116–117, on personal pronouns used reflexively.
% mdparagraph:P167:end

% mdparagraph:P168:begin
[10] Cook’s vocabulary, p.306 /\allowbreak{} PDF322, and \emph{The Day of Judgment}, p.136 /\allowbreak{} PDF152, printed line 6, for \textbf{ne} before the finite verb.
% mdparagraph:P168:end

% mdparagraph:P169:begin
[11] Cook, \emph{Trades and Occupations}, p.132 /\allowbreak{} PDF148, speaker headings before printed lines 2 and 13, for \textbf{sęcgan} introducing a direct utterance. The examples and both connected practice readings in this unit are new teaching material, not quotations from Cook.
% mdparagraph:P169:end

% mdparagraph:P170:begin
[12] Cook’s vocabulary, especially pp.290,292,294,296,299,305–306,312 /\allowbreak{} PDF306,308,310,312,315,321–322,328, checks the assigned lexical senses. Dictionary headword hyphens in prefixed verbs are joined for presentation.
% mdparagraph:P170:end

% mdparagraph:P171:begin
[13] Cook, Reader I, p.128 /\allowbreak{} PDF144, entire final paragraph, prose lines 1–7; only line marker 5 is printed. Source apparatus cues 1–13 are omitted. The original quantities, \textbf{ę}, punctuation and \textbf{for-ðon-ðe} are retained. For the opening passive see also §187, p.101 /\allowbreak{} PDF117.
% mdparagraph:P171:end

% mdparagraph:P172:begin
[14] Cook p.128 notes 3,5 and §§80,172.1, pp.48,97 /\allowbreak{} PDF64,113. The temporal accusative is explicitly noted by the source.
% mdparagraph:P172:end

% mdparagraph:P173:begin
[15] Cook §87(d), p.52 /\allowbreak{} PDF68, and §188, p.101 /\allowbreak{} PDF117.
% mdparagraph:P173:end

% mdparagraph:P174:begin
[16] Cook p.128 note 11 and §156(k), p.92 /\allowbreak{} PDF108, explicitly give the genitive after “cease”.
% mdparagraph:P174:end

% mdparagraph:P175:begin
[17] Cook p.128 note 12 gives Latin \textbf{creavit ut faceret} and “created to make”; §200, p.106 /\allowbreak{} PDF122, treats the inflected infinitive, called “gerund” there. The supplied source wording is \textbf{tō wyrceanne}, including its \textbf{e}.
% mdparagraph:P175:end

\endgroup
\end{document}
